\pagestyle{empty}
\tight
\begin{tblr}{width=\textwidth,colspec={|Q[c,m,1.9cm]|X[l,m]|},hlines,vlines,hspan=minimal,row{1}={bg=headblue},rowsep=1pt,colsep=3pt}
\SetCell{c,m}\hfil\logo\hfil & \textbf{POLITEKNIK NEGERI MALANG}\par\textbf{JURUSAN TEKNOLOGI INFORMASI}\par\textbf{PROGRAM STUDI: D4 TEKNIK INFORMATIKA} \\
\SetCell[c=2]{c} \textbf{RENCANA PEMBELAJARAN SEMESTER (RPS)} & \\
\end{tblr}
\begin{longtblr}{width=\textwidth,colspec={|Q[l,m,1.9cm]|X[0.85,l,m]|X[1.45,l,m]|X[0.75,l,m]|X[1.15,l,m]|},hlines,vlines,hspan=minimal,rowsep=1pt,colsep=2pt,row{1,3}={bg=headgray,font=\bfseries}}
MATA KULIAH & KODE & BOBOT (SKS)/jam & SEMESTER & TGL. PENYUSUNAN \\
Praktikum Dasar Pemrograman & RTI251007 & 3 SKS/6 jam & 1 & 7 Agustus 2025 \\
OTORISASI & Dosen Pengembang RPS & Koordinator RMK & \SetCell[c=2]{l} Ka PRODI &  \\
 & Mungki Astiningrum, S.T., M.Kom.\par Imam Fahrur Rozi, S.T., M.T.\par Vivi Nur Wijayaningrum, S.Kom., M.Kom.\par Mamluatul Hani'ah, S.Kom., M.Kom. & Imam Fahrur Rozi, S.T., M.T. & \SetCell[c=2]{l}{\vspace{8mm}\par Dr. Ely Setyo Astuti, S.T., M.T.} &  \\
\SetCell[r=11]{l} Capaian Pembelajaran (CP) & \SetCell[c=4]{l,bg=headgray,font=\bfseries} Capaian Pembelajaran Lulusan Program Studi (CPL-Prodi) &  &  &  \\
 & CPL07 & \SetCell[c=3]{l} Mampu menerapkan sistem komputasi cerdas dan pengolahan data untuk mendukung penyelesaian masalah komputasi. &  &  \\
 & CPL09 & \SetCell[c=3]{l} Mampu menjelaskan cara kerja sistem komputer dan menerapkan pendekatan komputasi untuk menyelesaikan masalah. &  &  \\
 & \SetCell[c=4]{l,bg=headgray,font=\bfseries} Capaian Pembelajaran mata kuliah (CPL-MK) &  &  &  \\
 & CPMK0704 & \SetCell[c=3]{l} Mahasiswa mampu menerapkan berbagai bahasa dan paradigma pemrograman untuk menyelesaikan masalah komputasi. &  &  \\
 & CPMK0903 & \SetCell[c=3]{l} Mahasiswa mampu menjelaskan prinsip logika komputasi serta menerapkan teknik dasar pemrograman dan algoritma untuk menganalisis dan menyelesaikan permasalahan. &  &  \\
 & \SetCell[c=4]{l,bg=headgray,font=\bfseries} Kemampuan akhir tiap tahapan belajar (Sub-CPMK) &  &  &  \\
 & SCPMK0903-00701 & \SetCell[c=3]{l} Mahasiswa mampu menjelaskan konsep program, bahasa pemrograman, compiler, dan debugging, melakukan instalasi tools pemrograman Java, serta memodelkan permasalahan sederhana menjadi algoritma (input, proses, output) [C2, P2]. &  &  \\
 & SCPMK0704-00702 & \SetCell[c=3]{l} Mahasiswa mampu menerapkan tipe data, variabel, operator, input-output, sequence, serta struktur pemilihan (if, if-else, else-if, switch-case, dan pemilihan bersarang) ke dalam program Java [C3, P3]. &  &  \\
 & SCPMK0704-00703 & \SetCell[c=3]{l} Mahasiswa mampu menerapkan struktur perulangan (for, while, do-while, perulangan bersarang) serta array 1 dimensi dan 2 dimensi untuk menyelesaikan studi kasus dalam program Java [C3, P3]. &  &  \\
 & SCPMK0704-00704 & \SetCell[c=3]{l} Mahasiswa mampu menerapkan fungsi, parameter, return, dan fungsi rekursif, serta membangun program Java utuh untuk menyelesaikan studi kasus dan mendemonstrasikannya [C3, P4]. &  &  \\
\SetCell[r=2]{l} Penilaian & \SetCell[c=4]{l} *Nilai CPMK didapat dari agregasi nilai SUB-CPMK yang dibebankan &  &  &  \\
 & \SetCell[c=4]{l}{\tiny\begin{tblr}{width=\linewidth,colspec={|Q[c,m,0.1\linewidth]|Q[c,m,0.16\linewidth]|X[l,m]|X[l,m]|X[l,m]|X[l,m]|X[l,m]|X[l,m]|Q[c,m,0.08\linewidth]|},hlines,vlines,hspan=minimal,rowsep=1pt,colsep=0.7pt,row{1,2}={bg=headgray,font=\bfseries}}
Kode CPMK & Kode SCPMK & \SetCell[c=6]{c} Bobot per Bentuk Penilaian & & & & & & Total Bobot \\
 & & Tugas Praktikum & Kuis 1 & UTS Praktik & Kuis 2 & PBL & UAS & \\
\SetCell[r=1]{c} CPMK0903 & SCPMK0903-00701 & 4\%: jobsheet instalasi Java dan studi kasus algoritma & 4\%: pemodelan algoritma & & & & & 8\% \\
\SetCell[r=3]{c} CPMK0704 & SCPMK0704-00702 & 6\%: jobsheet sequence dan pemilihan & 6\%: tipe data, IO, sequence & 10\%: implementasi pemilihan & & & & 22\% \\
 & SCPMK0704-00703 & 8\%: jobsheet perulangan dan array & & 10\%: implementasi perulangan & 10\%: perulangan bersarang dan array & & & 28\% \\
 & SCPMK0704-00704 & 4\%: jobsheet fungsi dan rekursif & & & & 8\%: proyek akhir & 30\%: demo project & 42\% \\
\SetCell[c=8]{r}\textbf{Total Per Penilaian} & & & & & & & & \textbf{100\%} \\
\end{tblr}} &  &  &  \\
Diskripsi Singkat Mata Kuliah & \SetCell[c=4]{l} Mata kuliah ini membahas pembuatan algoritma untuk menyelesaikan masalah dengan metode yang efisien hingga menghasilkan flowchart dan mentranslasikan teks algoritma ke dalam teks program bahasa pemrograman Java melalui praktik jobsheet terstruktur. &  &  &  \\
Bahan Kajian / Materi Pembelajaran & \SetCell[c=4]{l}\begin{enumerate}\item Konsep algoritma, representasi algoritma, translator, dan bahasa pemrograman\item Tipe data, variabel, konstanta, nilai, ekspresi, dan input-output\item Sequence, analisa kasus, pencabangan, dan perulangan\item Array dan fungsi/prosedur\end{enumerate} &  &  &  \\
\SetCell[r=2]{l} Pustaka & \SetCell[c=4]{l} \textbf{Utama:}\begin{enumerate}\item Sebesta, Robert. 2016. Concept of Programming Languages, Edisi Global. Addison Wesley.\item Sestoft, Peter. 2017. Programming Language Concepts. Springer.\item T. Henny Febriana Harumy. 2016. Belajar Dasar Algoritma dan Pemrograman C++. Deepublish.\end{enumerate} &  &  &  \\
 & \SetCell[c=4]{l} \textbf{Pendukung:} Rinaldi Munir. 2015. Algoritma dan Pemrograman. Penerbit Informatika; jobsheet praktikum dan materi LMS. &  &  &  \\
Media Pembelajaran & \SetCell[c=2]{l} \textbf{Software:} JDK, IDE, Microsoft Office, Adobe Reader, LMS. &  & \SetCell[c=2]{l} \textbf{Hardware:} PC/Laptop, LCD projector. &  \\
Nama Dosen Pengampu & \SetCell[c=4]{l}\begin{enumerate}\item Mungki Astiningrum, S.T., M.Kom.\item Imam Fahrur Rozi, S.T., M.T.\item Triana Fatmawati, S.T., M.T.\item Mamluatul Hani'ah, S.Kom., M.Kom.\item Rudy Ariyanto, S.T., M.Cs.\item Noprianto, S.Kom., M.Eng.\end{enumerate} &  &  &  \\
Matakuliah Syarat & \SetCell[c=4]{l} RTI251006 - Dasar Pemrograman &  &  &  \\
\end{longtblr}
\begin{longtblr}{width=\textwidth,colspec={|Q[c,m,0.06\textwidth]|X[1.35,l,m]|X[1.08,l,m]|X[1.16,l,m]|X[0.7,l,m]|X[1.24,l,m]|X[1.06,l,m]|X[0.98,l,m]|Q[c,m,0.05\textwidth]|},hlines,vlines,hspan=minimal,rowhead=2,row{1,2}={bg=headgreen,font=\bfseries},rowsep=1pt,colsep=1.5pt}
Mgg.\par Ke & Kemampuan Akhir Yang Direncanakan (Sub-CP-MK) & Bahan kajian (Materi Pembelajaran) & Bentuk dan Metode Pembelajaran & Estimasi Waktu & Pengalaman Belajar Mahasiswa & Kriteria \& Bentuk Penilaian & Indikator Penilaian & Bobot Penilaian (\%) \\
(1) & (2) & (3) & (4) & (5) & (6) & (7) & (8) & (9) \\
1 & SCPMK0903-00701 & Dasar pemrograman: konsep bahasa pemrograman, bahasa Java, compiler, instalasi tools pemrograman Java. & Modalitas: Luring/Daring. Metode: Problem Based Learning. Aktivitas: praktikum jobsheet di laboratorium. & 1 x 6 x 50' praktikum. & Mahasiswa melakukan instalasi tools Java sesuai jobsheet, melakukan compile dan debugging sintaks program, serta mengerjakan tugas praktikum. & Rubrik kriteria penilaian. Bentuk: jobsheet praktikum dan tugas mandiri. & Ketepatan instalasi tools pemrograman Java; keaktifan diskusi (tanya jawab). & 2 \\
2 & SCPMK0903-00701 & Studi kasus: pemodelan permasalahan dengan algoritma (input, proses, output). & Modalitas: Luring/Daring. Metode: Diskusi dan Problem Based Learning. Aktivitas: praktikum jobsheet. & 1 x 6 x 50' praktikum. & Mahasiswa melakukan langkah percobaan sesuai jobsheet dan membuat algoritma deskripsi berdasarkan studi kasus. & Rubrik kriteria penilaian. Bentuk: jobsheet praktikum dan tugas mandiri. & Ketepatan melaksanakan langkah percobaan; ketepatan menjawab pertanyaan pada jobsheet. & 2 \\
3 & SCPMK0704-00702 & Input-output, variabel, sequence, dan operator (aritmatika, gabungan, increment, decrement, relasional, logika, kondisional, bitwise, casting). & Modalitas: Luring/Daring. Metode: Diskusi. Aktivitas: praktikum jobsheet. & 1 x 6 x 50' praktikum. & Mahasiswa menggunakan tipe data, variabel, input-output, dan operator serta mentranslasikan permasalahan sequence ke bahasa Java. & Rubrik kriteria penilaian. Bentuk: jobsheet praktikum dan tugas mandiri. & Ketepatan menggunakan tipe data, variabel, input-output, operator; ketepatan kode program dan output. & 2 \\
4 & SCPMK0903-00701, SCPMK0704-00702 & Kuis 1: tes praktik materi minggu ke-1 sampai dengan 3. & Modalitas: Luring. Metode: tes praktik individu. & 1 x 6 x 50' praktikum. & Mahasiswa menyelesaikan soal tes praktik pemodelan algoritma, tipe data, input-output, dan sequence. & Rubrik kriteria penilaian. Bentuk: kuis (tes praktik). & Ketepatan pemodelan algoritma; ketepatan implementasi kode dan output. & 10 \\
5 & SCPMK0704-00702 & Pemilihan 1: if, if-else, else-if, dan switch-case. & Modalitas: Luring/Daring. Metode: Diskusi dan Problem Based Learning. Aktivitas: praktikum jobsheet. & 1 x 6 x 50' praktikum. & Mahasiswa menuliskan program Java dari flowchart kasus pemilihan dasar dan mengerjakan tugas praktikum. & Rubrik kriteria penilaian. Bentuk: jobsheet praktikum dan tugas mandiri. & Ketepatan menggunakan pemilihan 1; ketepatan kode program dan output sesuai tugas. & 2 \\
6 & SCPMK0704-00702 & Pemilihan 2: if bersarang. & Modalitas: Luring/Daring. Metode: Diskusi dan Problem Based Learning. Aktivitas: praktikum jobsheet. & 1 x 6 x 50' praktikum. & Mahasiswa menuliskan program Java dari flowchart kasus pemilihan bersarang (contoh: kasir sederhana). & Rubrik kriteria penilaian. Bentuk: jobsheet praktikum dan tugas mandiri. & Ketepatan menggunakan pemilihan bersarang; ketepatan menjawab pertanyaan jobsheet. & 2 \\
7 & SCPMK0704-00703 & Perulangan 1: for, while, do-while. & Modalitas: Luring/Daring. Metode: Diskusi. Aktivitas: praktikum jobsheet. & 1 x 6 x 50' praktikum. & Mahasiswa menuliskan program Java berdasarkan flowchart kasus perulangan bagian 1 dan memanfaatkan konsep perulangan. & Rubrik kriteria penilaian. Bentuk: jobsheet praktikum dan tugas mandiri. & Ketepatan menjelaskan perulangan bagian 1; ketepatan kode program dan output. & 2 \\
8 & SCPMK0704-00702, SCPMK0704-00703 & Ujian Tengah Semester: ujian praktik materi pemilihan dan perulangan. & Modalitas: Luring. Metode: ujian praktik individu. & 1 x 6 x 50' praktikum. & Mahasiswa menyelesaikan soal ujian praktik pembuatan program pemilihan dan perulangan. & Rubrik kriteria penilaian. Bentuk: UTS (ujian praktik). & Ketepatan analisis kasus; ketepatan implementasi pemilihan dan perulangan; kode berjalan sesuai output. & 20 \\
9 & SCPMK0704-00703 & Perulangan 2: perulangan bersarang (nested loop). & Modalitas: Luring/Daring. Metode: Diskusi. Aktivitas: praktikum jobsheet. & 1 x 6 x 50' praktikum. & Mahasiswa menuliskan program Java berdasarkan flowchart kasus perulangan bersarang. & Rubrik kriteria penilaian. Bentuk: jobsheet praktikum dan tugas mandiri. & Ketepatan menjelaskan perulangan bersarang; ketepatan memodelkan studi kasus ke Java. & 2 \\
10 & SCPMK0704-00703 & Array 1: array 1 dimensi, pengaksesan elemen, pengayaan searching dan sorting. & Modalitas: Luring/Daring. Metode: Diskusi. Aktivitas: praktikum jobsheet. & 1 x 6 x 50' praktikum. & Mahasiswa menerapkan array 1 dimensi menggunakan bahasa Java sesuai jobsheet dan tugas praktikum. & Rubrik kriteria penilaian. Bentuk: jobsheet praktikum dan tugas mandiri. & Ketepatan menggunakan array 1 dimensi; ketepatan kode program dan output. & 2 \\
11 & SCPMK0704-00703 & Array 2: array 2 dimensi, pengayaan operasi matriks. & Modalitas: Luring/Daring. Metode: Diskusi dan Problem Based Learning. Aktivitas: praktikum jobsheet. & 1 x 6 x 50' praktikum. & Mahasiswa menerapkan konsep array 2 dimensi pada studi kasus yang diberikan (contoh: pencarian data usia termuda). & Rubrik kriteria penilaian. Bentuk: jobsheet praktikum dan tugas mandiri. & Ketepatan menggunakan array 2 dimensi; ketepatan kode program dan output. & 2 \\
12 & SCPMK0704-00703 & Kuis 2: tes praktik perulangan bersarang dan array. & Modalitas: Luring. Metode: tes praktik individu. & 1 x 6 x 50' praktikum. & Mahasiswa menyelesaikan soal tes praktik perulangan bersarang, array 1 dimensi, dan array 2 dimensi. & Rubrik kriteria penilaian. Bentuk: kuis (tes praktik). & Ketepatan implementasi perulangan bersarang; ketepatan implementasi array. & 10 \\
13 & SCPMK0704-00704 & Fungsi 1: tipe data fungsi, parameter/argumen, return, dan pemanggilan fungsi. & Modalitas: Luring/Daring. Metode: Diskusi. Aktivitas: praktikum jobsheet. & 1 x 6 x 50' praktikum. & Mahasiswa menggunakan fungsi pada bahasa Java sesuai jobsheet dan mengerjakan tugas praktikum. & Rubrik kriteria penilaian. Bentuk: jobsheet praktikum dan tugas mandiri. & Ketepatan menggunakan fungsi; ketepatan kode program dan output sesuai tugas. & 2 \\
14 & SCPMK0704-00704 & Fungsi 2: fungsi rekursif dan pengayaan kasus fungsi. & Modalitas: Luring/Daring. Metode: Diskusi. Aktivitas: praktikum jobsheet. & 1 x 6 x 50' praktikum. & Mahasiswa menggunakan fungsi rekursif pada bahasa Java sesuai jobsheet dan mengerjakan tugas praktikum. & Rubrik kriteria penilaian. Bentuk: jobsheet praktikum dan tugas mandiri. & Ketepatan menggunakan fungsi rekursif; keaktifan diskusi; ketepatan kode dan output. & 2 \\
15 & SCPMK0704-00704 & Proyek akhir: pembuatan program penyelesaian permasalahan (materi pertemuan 1--14). & Modalitas: Luring/Daring. Metode: Project Based Learning, praktek mandiri. & 1 x 6 x 50' praktikum. & Mahasiswa membuat flowchart dan program untuk menyelesaikan studi kasus sesuai topik yang ditentukan. & Rubrik kriteria penilaian. Bentuk: PBL proyek akhir. & Ketepatan rancangan flowchart proyek; ketepatan kode program sesuai studi kasus. & 4 \\
16 & SCPMK0704-00704 & Penyelesaian proyek akhir dan demo program (materi pertemuan 1--14). & Modalitas: Luring/Daring. Metode: Project Based Learning, praktek mandiri, demo program. & 1 x 6 x 50' praktikum. & Mahasiswa menyelesaikan program proyek akhir, melakukan latihan, dan mendemonstrasikan program. & Rubrik kriteria penilaian. Bentuk: PBL proyek akhir. & Keaktifan diskusi; ketepatan kode program dan output sesuai latihan; kesiapan demo. & 4 \\
17 & SCPMK0704-00704 & Ujian Akhir Semester: demo project. & Modalitas: Luring. Metode: ujian praktik dan demo project individu. & 1 x 6 x 50' praktikum. & Mahasiswa mendemonstrasikan program proyek akhir dan menjawab pertanyaan penguji terkait kode program. & Rubrik kriteria penilaian. Bentuk: UAS (demo project). & Kelengkapan fungsionalitas program; kualitas kode; kualitas demo dan penjelasan program. & 30 \\
\end{longtblr}
